% Options for packages loaded elsewhere
\PassOptionsToPackage{unicode}{hyperref}
\PassOptionsToPackage{hyphens}{url}
\PassOptionsToPackage{dvipsnames,svgnames,x11names}{xcolor}
\documentclass[
  10pt,
  fontset=fandol]{ctexart}
\usepackage{xcolor}
\usepackage[margin=16mm]{geometry}
\usepackage{amsmath,amssymb}
\setcounter{secnumdepth}{-\maxdimen} % remove section numbering
\usepackage{iftex}
\ifPDFTeX
  \usepackage[T1]{fontenc}
  \usepackage[utf8]{inputenc}
  \usepackage{textcomp} % provide euro and other symbols
\else % if luatex or xetex
  \usepackage{unicode-math} % this also loads fontspec
  \defaultfontfeatures{Scale=MatchLowercase}
  \defaultfontfeatures[\rmfamily]{Ligatures=TeX,Scale=1}
\fi
\usepackage{lmodern}
\ifPDFTeX\else
  % xetex/luatex font selection
\fi
% Use upquote if available, for straight quotes in verbatim environments
\IfFileExists{upquote.sty}{\usepackage{upquote}}{}
\IfFileExists{microtype.sty}{% use microtype if available
  \usepackage[]{microtype}
  \UseMicrotypeSet[protrusion]{basicmath} % disable protrusion for tt fonts
}{}
\makeatletter
\@ifundefined{KOMAClassName}{% if non-KOMA class
  \IfFileExists{parskip.sty}{%
    \usepackage{parskip}
  }{% else
    \setlength{\parindent}{0pt}
    \setlength{\parskip}{6pt plus 2pt minus 1pt}}
}{% if KOMA class
  \KOMAoptions{parskip=half}}
\makeatother
\setlength{\emergencystretch}{3em} % prevent overfull lines
\providecommand{\tightlist}{%
  \setlength{\itemsep}{0pt}\setlength{\parskip}{0pt}}
\usepackage{amsmath,amssymb,mathtools,needspace}
\xeCJKDeclareCharClass{CJK}{"2032,"0400->"04FF,"2160->"216B,"2460->"2473,"25A0,"25C7,"25CB,"25CF}
\IfFontExistsTF{Arial Unicode MS}{\xeCJKsetup{AutoFallBack=true}\setCJKfallbackfamilyfont{\CJKrmdefault}{Arial Unicode MS}}{}
\setcounter{secnumdepth}{0}
\setlength{\emergencystretch}{2em}
\allowdisplaybreaks
\usepackage{bookmark}
\IfFileExists{xurl.sty}{\usepackage{xurl}}{} % add URL line breaks if available
\urlstyle{same}
\hypersetup{
  pdftitle={三转角方程},
  colorlinks=true,
  linkcolor={Maroon},
  filecolor={Maroon},
  citecolor={Blue},
  urlcolor={Blue},
  pdfcreator={LaTeX via pandoc}}

\title{三转角方程}
\author{}
\date{}

\begin{document}
\maketitle

\subsubsection{2.8.2 三转角方程}\label{ux4e09ux8f6cux89d2ux65b9ux7a0b}

现在构造满足条件（2.8.1）及加上相应边界条件的三次样条函数 \(S(x)\) 的表达式。若假定 \(S'(x)\) 在节点 \(x_j\) 处的值为 \(S'(x_j)=m_j\)（\(j=0,1,\cdots,n\)），再由式（2.8.1），则由分段三次 Hermite 插值式（2.7.6）可得

\[S(x)=\sum_{j=0}^n[y_j\alpha_j(x)+m_j\beta_j(x)],\tag{2.8.6}\]

其中 \(\alpha_j(x),\beta_j(x)\) 是插值基函数，由式（2.7.7）、式（2.7.8）表示。显然，表达式（2.8.6）中 \(S(x)\) 及 \(S'(x)\) 在整个区间 \([a,b]\) 上连续，且满足式（2.8.1）；然而在式（2.8.6）中 \(m_j\)（\(j=0,1,\cdots,n\)）是未知的，它可利用式（2.8.2）

\[S''(x_j-0)=S''(x_j+0)\quad(j=1,\cdots,n-1)\]

及边界条件（2.8.3）（也可以是其他边界条件）来确定。为了求出 \(m_j\)，下面考虑 \(S(x)\) 在 \([x_j,x_{j+1}]\) 上的表达式

\[\begin{aligned}
S(x)={}&\frac{(x-x_{j+1})^2[h_j+2(x-x_j)]}{h_j^3}y_j+\frac{(x-x_j)^2[h_j+2(x_{j+1}-x)]}{h_j^3}y_{j+1}\\
&+\frac{(x-x_{j+1})^2(x-x_j)}{h_j^2}m_j+\frac{(x-x_j)^2(x-x_{j+1})}{h_j^2}m_{j+1},
\end{aligned}\tag{2.8.7}\]

这里 \(h_j=x_{j+1}-x_j\)。对 \(S(x)\) 求二阶导数，得

\[\begin{aligned}
S''(x)={}&\frac{6x-2x_j-4x_{j+1}}{h_j^2}m_j+\frac{6x-4x_j-2x_{j+1}}{h_j^2}m_{j+1}\\
&+\frac{6(x_j+x_{j+1}-2x)}{h_j^3}(y_{j+1}-y_j),
\end{aligned}\]

于是

\[S''(x_j+0)=-\frac4{h_j}m_j-\frac2{h_j}m_{j+1}+\frac6{h_j^2}(y_{j+1}-y_j).\]

同理，可得 \(S''(x)\) 在区间 \([x_{j-1},x_j]\) 上的表达式

\[\begin{aligned}
S''(x)={}&\frac{6x-2x_{j-1}-4x_j}{h_{j-1}^2}m_{j-1}+\frac{6x-4x_{j-1}-2x_j}{h_{j-1}^2}m_j\\
&+\frac{6(x_{j-1}+x_j-2x)}{h_{j-1}^2}(y_j-y_{j-1})
\end{aligned}\]

\begin{quote}
续页：左侧二阶导数在下一页继续。
\end{quote}

\begin{quote}
\textbf{校注（agent 补充；原书疑误）}：本页末式第三项的分母在原扫描中为 \(h_{j-1}^2\)，上面原样保留，参见\href{../assets/ch02b/pdf-050-second-derivative-detail.png}{原式放大裁图}。由同页前一个区间的 \(S''(x)\) 公式换下标，第三项分母应为 \(h_{j-1}^3\)；这样代入 \(x=x_j\) 才得到下一页的 \(-6(y_j-y_{j-1})/h_{j-1}^2\)。这是原书幂次疑误；OCR 还漏掉了多个 \(j-1\) 下标，已按图补回。
\end{quote}

\href{../../source-images/pdf-051.jpeg}{查看原页}

\begin{quote}
本页接上页 2.8.2 正文。
\end{quote}

及

\[S''(x_j-0)=\frac2{h_{j-1}}m_{j-1}+\frac4{h_{j-1}}m_j-\frac6{h_{j-1}^2}(y_j-y_{j-1}).\]

由条件 \(S''(x_j+0)=S''(x_j-0)\)（\(j=1,2,\cdots,n-1\)），可得

\[\begin{aligned}
&\frac1{h_{j-1}}m_{j-1}+2\left(\frac1{h_{j-1}}+\frac1{h_j}\right)m_j+\frac1{h_j}m_{j+1}\\
&\qquad=3\left(\frac{y_{j+1}-y_j}{h_j^2}+\frac{y_j-y_{j-1}}{h_{j-1}^2}\right)\quad(j=1,2,\cdots,n-1),
\end{aligned}\tag{2.8.8}\]

用 \(\frac1{h_{j-1}}+\frac1{h_j}\) 除全式，并注意 \(y_j=f_j,\frac{y_{j+1}-y_j}{h_j}=f[x_j,x_{j+1}]\)，方程（2.8.8）可简化为

\[\lambda_jm_{j-1}+2m_j+\mu_jm_{j+1}=g_j\quad(j=1,2,\cdots,n-1),\tag{2.8.9}\]

其中

\[\lambda_j=\frac{h_j}{h_{j-1}+h_j},\quad\mu_j=\frac{h_{j-1}}{h_{j-1}+h_j}\quad(j=1,2,\cdots,n-1),\tag{2.8.10}\]

\[g_j=3(\lambda_jf[x_{j-1},x_j]+\mu_jf[x_j,x_{j+1}])\quad(j=1,2,\cdots,n-1).\tag{2.8.11}\]

方程（2.8.9）是关于未知数 \(m_0,m_1,\cdots,m_n\) 的 \(n-1\) 个方程，若加上已知条件（2.8.3）：\(m_0=f'_0\) 则 \(m_n=f'_n\)，那么方程（2.8.9）为只含 \(m_1,\cdots,m_{n-1}\) 的 \(n-1\) 个方程，写成矩阵形式便是

\[\begin{pmatrix}
2&\mu_1&0&\cdots&\cdots&0\\
\lambda_2&2&\mu_2&\ddots&&\vdots\\
0&\lambda_3&2&\mu_3&\ddots&\vdots\\
\vdots&\ddots&\ddots&\ddots&\ddots&0\\
0&&\ddots&\lambda_{n-2}&2&\mu_{n-2}\\
0&\cdots&\cdots&0&\lambda_{n-1}&2
\end{pmatrix}
\begin{pmatrix}m_1\\m_2\\m_3\\\vdots\\m_{n-2}\\m_{n-1}\end{pmatrix}
=\begin{pmatrix}g_1-\lambda_1f'_0\\g_2\\g_3\\\vdots\\g_{n-2}\\g_{n-1}-\mu_{n-1}f'_n\end{pmatrix}.\tag{2.8.12}\]

如果边界条件为式（2.8.4），则

\[\begin{cases}
\displaystyle2m_0+m_1=3f[x_0,x_1]-\frac{h_0}2f''_0=g_0,\\[6pt]
\displaystyle m_{n-1}+2m_n=3f[x_{n-1},x_n]+\frac{h_{n-1}}2f''_n=g_n.
\end{cases}\tag{2.8.13}\]

如果边界条件为式（2.8.4）′，则

\[\begin{cases}
2m_0+m_1=3f[x_0,x_1]=g_0,\\
m_{n-1}+2m_n=3f[x_{n-1},x_n]=g_n.
\end{cases}\tag{2.8.13'}\]

于是，式（2.8.9）与式（2.8.13）或式（2.8.13）′合并后用矩阵形式表示为

\[\begin{pmatrix}
2&1&0&\cdots&\cdots&0\\
\lambda_1&2&\mu_1&\ddots&&\vdots\\
0&\lambda_2&2&\mu_2&\ddots&\vdots\\
\vdots&\ddots&\ddots&\ddots&\ddots&0\\
\vdots&&\ddots&\lambda_{n-1}&2&\mu_{n-1}\\
0&\cdots&\cdots&0&1&2
\end{pmatrix}
\begin{pmatrix}m_0\\m_1\\m_2\\\vdots\\m_{n-1}\\m_n\end{pmatrix}
=\begin{pmatrix}g_0\\g_1\\g_2\\\vdots\\g_{n-1}\\g_n\end{pmatrix}.\tag{2.8.14}\]

\href{../../source-images/pdf-052.jpeg}{查看原页}

\begin{quote}
本页接上页 2.8.2 正文。
\end{quote}

如果边界条件为周期性条件式（2.8.5），则

\[m_0=m_n,\]

\[\frac1{h_0}m_1+\frac1{h_{n-1}}m_{n-1}+2\left(\frac1{h_0}+\frac1{h_{n-1}}\right)m_n=\frac3{h_0}f[x_0,x_1]+\frac3{h_{n-1}}f[x_{n-1},x_n],\]

化简为

\[\mu_nm_1+\lambda_nm_{n-1}+2m_n=g_n,\]

\[\mu_n=\frac{h_{n-1}}{h_0+h_{n-1}},\quad\lambda_n=\frac{h_{n-1}}{h_0+h_{n-1}},\]

\[g_n=3(\mu_nf[x_0,x_1]+\lambda_nf[x_{n-1},x_n]).\]

与式（2.8.9）合并，用矩阵形式表示为

\[\begin{pmatrix}
2&\mu_1&0&\cdots&0\\
\lambda_2&2&\mu_2&\ddots&\vdots\\
0&\lambda_3&\ddots&\ddots&0\\
\vdots&\ddots&\ddots&2&\mu_{n-1}\\
0&\cdots&0&\lambda_n&2
\end{pmatrix}
\begin{pmatrix}m_1\\m_2\\\vdots\\m_{n-1}\\m_n\end{pmatrix}
=\begin{pmatrix}g_1\\g_2\\\vdots\\g_{n-1}\\g_n\end{pmatrix}.\tag{2.8.15}\]

这里得到的方程组（2.8.12）、（2.8.14）及式（2.8.15）中，每个方程都联系三个 \(m_j\)，\(m_j\) 在力学上解释为细梁在 \(x_j\) 截面处的转角，故称之为\textbf{三转角方程}。这些方程系数矩阵对角元素均为 2，非对角元素 \(\mu_j+\lambda_j=1\)，故系数矩阵具有强对角优势，方程组（2.8.12）、（2.8.14）及（2.8.15）都有唯一解，可用追赶法求得解 \(m_j\)（\(j=0,1,\cdots,n\)），从而得到 \(S(x)\)。关于线性方程组的解法见本书第 7 章。

\begin{center}\rule{0.5\linewidth}{0.5pt}\end{center}

\subsection{相关原页校注（agent 补充）}\label{ux76f8ux5173ux539fux9875ux6821ux6ce8agent-ux8865ux5145}

以下校注来自本节覆盖的原页，可能也涉及同页相邻小节；原文照录与校正说明分开保留。

\subsubsection{原书 PDF 页 52 的校注}\label{ux539fux4e66-pdf-ux9875-52-ux7684ux6821ux6ce8}

\href{../pages/pdf-052.md}{完整原页转录} · \href{../../source-images/pdf-052.jpeg}{原页图像}

校注记录：ch02b-E004, ch02b-E005。

\begin{quote}
\textbf{校注（agent 补充；原书疑误）}：本页 \(\lambda_n\) 原式与 \(\mu_n\) 同为 \(h_{n-1}/(h_0+h_{n-1})\)，已保留。由上一行未归一化方程，\(m_{n-1}\) 的系数应为 \(\lambda_n=h_0/(h_0+h_{n-1})\)。此外，（2.8.15）的右上角、左下角在原书均印为 \(0\)；由 \(m_0=m_n\) 及本页的周期方程，应分别为 \(\lambda_1\)、\(\mu_n\)。对应系统是有两个角元素的循环三对角系统，不能不作处理便直接套用普通三对角追赶法。参见\href{../assets/ch02b/pdf-052-periodic-system-detail.png}{原文和矩阵放大裁图}。
\end{quote}

\subsection{本节覆盖原页的校注记录}\label{ux672cux8282ux8986ux76d6ux539fux9875ux7684ux6821ux6ce8ux8bb0ux5f55}

下列记录按原页关联，含同页相邻小节的校注；计算时同时读取上文校注和完整原页。

\begin{itemize}
\tightlist
\item
  \texttt{ch02b-E003}：原书 PDF 页 50
\item
  \texttt{ch02b-E004}：原书 PDF 页 52
\item
  \texttt{ch02b-E005}：原书 PDF 页 52
\end{itemize}

\end{document}
